\paragraph{The two mass-twelve rows of relaxed cost $2604$.}
Let $G_-$ have mass eighty and $G_+$ have mass ninety-six,
in either orientation of the closed table. The two rows now considered have
\[
 W(4)=12,\qquad \delta_- =\psi+2R,\quad U_-=D+2,
 \qquad \delta_+=\psi,\quad U_+(4)=12,
\]
where $\psi=D(D-3)$ and $R$ is either the constant two or $T_k$,
with $k\ge1$. Every allowed auxiliary increment is included.
In particular $\delta_-(0)=4$ or $-4$.

We first record the divisibility interface used in both orientations.
Write $W=D+2B$ with $B\in\mathbb Z[D]$. Since $W(4)=12$,
$B(4)=4$; therefore $B(-2)$ is even, and
$W(-2)\equiv2\pmod4$. Thus $W$ is coprime to $D+2$
over $\mathbb Q[D]$. From the minus-cell identity
\[
                 (D+2)G_+=W(G_--D\delta_-)
\]
it follows that $W$ divides $G_+$ in $\mathbb Z[D]$.
Here integrality follows either by Gauss's lemma and content,
or directly as follows. Write $W=cW_0$, with $W_0$ primitive.
The rational coprimality gives $G_+=W_0Q$ with $Q$ integral.
After cancellation, $c$ divides every coefficient of $(D+2)Q$;
the product-content identity and primitivity of $D+2$ imply
that $c$ divides every coefficient of $Q$.

If $W(0)\ne0$, the preceding mass-twelve classification gives
$W(0)=\pm8$ and $U_+(0)\in\{-8,0,8\}$.
The two cell identities at zero force both $G_i(0)=0$:
their coefficient determinant is
$W(0)^2-U_-(0)U_+(0)=64-2U_+(0)\ne0$.
Write $G_\pm=DP_\pm$ integrally. Dividing the identities by
$D$ and evaluating at zero yields
\[
 wP_-(0)-2P_+(0)=w\delta_-(0),\qquad
 wP_+(0)-uP_-(0)=0,
 \quad w=W(0),\ u=U_+(0).
\]
Consequently
\[
 P_-(0)=\frac{64\delta_-(0)}{64-2u}.
\]
For $u=8$ or $-8$ this is respectively
$4\delta_-(0)/3$ or $4\delta_-(0)/5$, neither an integer.
Hence $u=0$ and $P_+(0)=0$, so $D^2$ divides $G_+$.
On the other hand $G_+=WQ$ with $Q(4)=96/12=8$.
Since $W(0)\ne0$, $D^2$ divides $Q$ integrally.
Evaluation at four would make sixteen divide eight, a contradiction.

It remains to treat $W(0)=0$, so $W=DH$ with $H$ as in
the preceding mass-twelve gluing lemma. If $H$ is nonconstant,
it is primitive irreducible with leading coefficient two,
$H(4)=3$, and odd $H(0)\ne1,-1$.
It cannot divide $G_-$, since three does not divide eighty.
The plus-cell identity therefore forces $H\mid U_+$.
Divisibility at zero, together with $U_+(0)\in\{-8,0,8\}$,
forces $U_+(0)=0$. Write $U_+=DJ$.
The same classification makes $J$ either three or a primitive
irreducible polynomial of leading coefficient two and value three
at four. Since $H\mid J$, the constant case is impossible and
$J=H$. Thus $U_+=W$.
If $H=3$, reduce the plus-cell identity modulo three.
The reduction of $G_-$ is nonzero because its mass is eighty,
so all coefficients of $U_+$ are divisible by three.
Its odd coefficients at $q,q^{-1}$ are at least three each,
and its even central coefficient at least six. These exhaust
mass twelve, whence again $U_+=3D=W$.

Cancelling $W$ in the plus-cell identity gives
$G_+=G_-+D\psi$. The other cell identity becomes
\[
 (W-D-2)G_-=D\bigl(W\delta_-+(D+2)\psi\bigr).
\]
Its right side is divisible by $D^2$, whereas the first factor
on the left has value $-2$ at zero. Therefore $D^2\mid G_-$.
If $G_-=G_1$, its quotient by $D^2$ has every coefficient even
but evaluation $80/16=5$ at four, a contradiction.
If $G_-=G_2$, the necessary congruence of NN-G030 is
\[
 G_--D\delta_-\equiv G_+\pmod4.
\]
Substituting $G_+=G_-+D\psi$ gives
$2D(\psi+R)\equiv0\pmod4$, so $R\equiv\psi\pmod2$.
This is impossible. The constant two reduces to zero; if
$R=T_k$, its monic degree is $k$, so equality with the degree-two
polynomial $\psi$ would require $k=2$. But
$T_2\equiv D^2$ and $\psi\equiv D^2+D$ modulo two.
Both orientations are excluded for arbitrary grading support.
\hfill$\square$

The objects and scope are those of NN-G036: actual effective graded
sign representations composing in the symmetric closed table.
Restriction at $q=-1$, ordinary dimension at $q=1$, and integral
character divisibility are operations in the torus representation
ring of SGA~3, I, 4.7.2--4.7.3. No quotient is assumed effective,
and no statement about extra-position tables is inferred.
